% ============================================================
%  VidSeek – Graduation Project Report
%  Cairo University, Faculty of Engineering
%  Department of Computer Engineering
% ============================================================
\documentclass[12pt, a4paper]{report}

% ── Packages ─────────────────────────────────────────────────────────────────
\usepackage[utf8]{inputenc}
\usepackage[T1]{fontenc}
\usepackage{geometry}
\usepackage{graphicx}
\usepackage{amsmath}
\usepackage{amssymb}
\usepackage{booktabs}
\usepackage{longtable}
\usepackage{array}
\usepackage{enumitem}
\usepackage{hyperref}
\usepackage{fancyhdr}
\usepackage{titlesec}
\usepackage{setspace}
\usepackage{parskip}
\usepackage{caption}
\usepackage{subcaption}
\usepackage{float}
\usepackage{xcolor}
\usepackage{listings}
\usepackage{algorithm}
\usepackage{algpseudocode}
\usepackage{tikz}
\usepackage{multirow}
\usepackage{lmodern}
\usepackage{microtype}
\usepackage[acronym]{glossaries}
\usepackage{pdfpages}
\usepackage{afterpage}

% ── Page geometry ─────────────────────────────────────────────────────────────
\geometry{
  top    = 2.5cm,
  bottom = 2.5cm,
  left   = 3.0cm,
  right  = 2.5cm,
}

% ── Line spacing ──────────────────────────────────────────────────────────────
\onehalfspacing

% ── Header / footer ──────────────────────────────────────────────────────────
\pagestyle{fancy}
\fancyhf{}
\fancyhead[L]{\small VidSeek \the\year}
\fancyhead[R]{}
\fancyfoot[R]{\thepage\ | Page}
\renewcommand{\headrulewidth}{0.4pt}

% ── Chapter title style ──────────────────────────────────────────────────────
\titleformat{\chapter}[display]
  {\normalfont\Large\bfseries}
  {\chaptername\ \thechapter}{10pt}{\LARGE}
\titlespacing*{\chapter}{0pt}{30pt}{20pt}

% ── Hyperref colours ─────────────────────────────────────────────────────────
\hypersetup{
  colorlinks = true,
  linkcolor  = black,
  urlcolor   = blue,
  citecolor  = black,
}

% ── Listings style ───────────────────────────────────────────────────────────
\lstset{
  basicstyle   = \ttfamily\small,
  breaklines   = true,
  frame        = single,
  numbers      = left,
  numberstyle  = \tiny\color{gray},
}

% =============================================================
\begin{document}

% ── Title Page ───────────────────────────────────────────────────────────────
\begin{titlepage}
  \centering

  \vspace*{0.5cm}
  {\Large\bfseries Cairo University\\[0.3em]
   Faculty of Engineering\\[0.3em]
   Department of Computer Engineering}

  \vspace{2.5cm}
  {\Huge\bfseries VidSeek}

  \vspace{0.8cm}
  {\large\itshape A Multi-Modal Video Search and Retrieval System}

  \vspace{2.5cm}
  {\large A Graduation Project Report Submitted\\[0.4em]
   to\\[0.4em]
   Faculty of Engineering, Cairo University\\[0.4em]
   in Partial Fulfillment of the requirements of the degree\\[0.4em]
   of\\[0.4em]
   Bachelor of Science in Computer Engineering.}

  \vspace{2.5cm}
  {\large Presented by}

  \vspace{0.5cm}
  {\large
    Amr Mahmoud\\[0.3em]
    % Add remaining team members here
  }

  \vspace{2cm}
  {\large Supervised by}

  \vspace{0.5cm}
  {\large Prof. [Supervisor Name]}

  \vfill
  {\large June 2026}

  \vspace{1cm}
  {\small All rights reserved. This report may not be reproduced in whole or in part,
  by photocopying or other means, without the permission of the authors/department.}
\end{titlepage}

% ── Front matter ─────────────────────────────────────────────────────────────
\pagenumbering{roman}

% ─── Abstract ────────────────────────────────────────────────────────────────
\chapter*{Abstract}
\addcontentsline{toc}{chapter}{Abstract}
\thispagestyle{fancy}

Video content on the internet is growing at an enormous rate, but searching
inside a video remains a hard problem. Unlike text documents, videos cannot
be searched with a simple keyword query because their meaning is spread
across speech, on-screen text, detected objects, and visual relationships
between those objects.

This project introduces VidSeek, a multi-modal video search and retrieval
system. VidSeek indexes videos by processing them through four parallel
pipelines: Optical Character Recognition (OCR) extracts text shown on
screen, Automatic Speech Recognition (ASR) transcribes spoken audio,
object detection locates named objects in each frame, and Visual Relationship
Detection (VRD) describes how those objects relate to each other. All
extracted information is converted into dense vector embeddings and stored
in a custom two-level Inverted File (IVF) index together with a relational
database.

At query time, a user types a natural-language question. The system
encodes it into the same embedding space, searches the IVF index, and
returns a ranked list of video segments - with timestamps - that are most
likely to contain the answer.

The result is a system that lets users search any uploaded video as easily
as they would search a text document.

\newpage

% ─── Acknowledgment ──────────────────────────────────────────────────────────
\chapter*{Acknowledgment}
\addcontentsline{toc}{chapter}{Acknowledgment}
\thispagestyle{fancy}

(Acknowledgment text to be added here.)

\newpage

% ─── Table of Contents ───────────────────────────────────────────────────────
\tableofcontents
\newpage

% ─── List of Figures ─────────────────────────────────────────────────────────
\listoffigures
\addcontentsline{toc}{chapter}{List of Figures}
\newpage

% ─── List of Tables ──────────────────────────────────────────────────────────
\listoftables
\addcontentsline{toc}{chapter}{List of Tables}
\newpage

% ─── List of Abbreviations ───────────────────────────────────────────────────
\chapter*{List of Abbreviations}
\addcontentsline{toc}{chapter}{List of Abbreviations}
\thispagestyle{fancy}

\begin{longtable}{>{\bfseries}p{3cm} p{10cm}}
  \toprule
  \textbf{Abbreviation} & \textbf{Meaning} \\
  \midrule
  \endfirsthead
  \toprule
  \textbf{Abbreviation} & \textbf{Meaning} \\
  \midrule
  \endhead
  ASR    & Automatic Speech Recognition \\
  CNN    & Convolutional Neural Network \\
  FAISS  & Facebook AI Similarity Search \\
  FPN    & Feature Pyramid Network \\
  GPU    & Graphics Processing Unit \\
  HOG    & Histogram of Oriented Gradients \\
  IVF    & Inverted File Index \\
  KG     & Knowledge Graph \\
  LLM    & Large Language Model \\
  MSER   & Maximally Stable Extremal Regions \\
  NLP    & Natural Language Processing \\
  OCR    & Optical Character Recognition \\
  PCA    & Principal Component Analysis \\
  ReLU   & Rectified Linear Unit \\
  RoI    & Region of Interest \\
  RPN    & Region Proposal Network \\
  SQL    & Structured Query Language \\
  SVM    & Support Vector Machine \\
  TF-IDF & Term Frequency–Inverse Document Frequency \\
  VRD    & Visual Relationship Detection \\
  \bottomrule
\end{longtable}

\newpage

% ─── List of Symbols ─────────────────────────────────────────────────────────
\chapter*{List of Symbols}
\addcontentsline{toc}{chapter}{List of Symbols}
\thispagestyle{fancy}

\begin{longtable}{>{\bfseries}p{3cm} p{10cm}}
  \toprule
  \textbf{Symbol} & \textbf{Meaning} \\
  \midrule
  \endfirsthead
  \toprule
  \textbf{Symbol} & \textbf{Meaning} \\
  \midrule
  \endhead
  $d$          & Embedding dimension \\
  $\mathbf{q}$ & Query vector, $\mathbf{q} \in \mathbb{R}^d$ \\
  $\mathbf{v}$ & Document vector, $\mathbf{v} \in \mathbb{R}^d$ \\
  $k$          & Number of nearest neighbours to return \\
  $N$          & Total number of indexed vectors \\
  $C_1$        & Number of level-1 IVF clusters \\
  $C_2$        & Number of level-2 IVF sub-clusters per level-1 cluster \\
  $p_1$        & Number of level-1 clusters probed at query time \\
  $p_2$        & Number of level-2 sub-clusters probed at query time \\
  $\delta$     & Threshold distance for creating a new cluster centroid \\
  $\cos(\cdot,\cdot)$ & Cosine similarity between two vectors \\
  $t_s, t_e$  & Start and end timestamps of a video segment (seconds) \\
  $F$          & Set of video frames extracted from a scene \\
  $S$          & Set of scene segments detected in a video \\
  \bottomrule
\end{longtable}

\newpage

% ─── Contacts ────────────────────────────────────────────────────────────────
\chapter*{Contacts}
\addcontentsline{toc}{chapter}{Contacts}
\thispagestyle{fancy}

\textbf{Team Members}

\vspace{0.5em}
\begin{tabular}{lll}
  \toprule
  \textbf{Name} & \textbf{Email} & \textbf{Phone Number} \\
  \midrule
  Amr Mahmoud & amrmgoma6@gmail.com & -- \\
  % Add remaining team members here
  \bottomrule
\end{tabular}

\vspace{2em}
\textbf{Supervisor}

\vspace{0.5em}
\begin{tabular}{lll}
  \toprule
  \textbf{Name} & \textbf{Email} & \textbf{Number} \\
  \midrule
  Prof. [Supervisor Name] & [supervisor@email] & -- \\
  \bottomrule
\end{tabular}

\newpage

% =============================================================
%  Main chapters
% =============================================================
\pagenumbering{arabic}
\setcounter{page}{1}

% ─────────────────────────────────────────────────────────────────────────────
\chapter{Introduction}
% ─────────────────────────────────────────────────────────────────────────────

\section{Motivation and Justification}

Video is now the most common format for sharing information on the internet.
Platforms such as YouTube report more than 500 hours of new video uploaded
every minute, and enterprise organisations store thousands of hours of
lecture recordings, training sessions, news broadcasts, and security footage.
Despite this enormous volume, video is fundamentally hard to search. When
a user wants to find a specific moment in a video they must either watch
the entire recording or rely on manually written captions, which are not
always available and rarely complete.

The difficulty comes from the nature of the content itself. A video carries
information across at least four distinct channels at the same time: the
spoken words of people on screen, the text that appears as titles or subtitles,
the objects that are visible in each frame, and the spatial relationships
between those objects. Current general-purpose search engines can only
index whatever short text description was attached to the video at upload
time. They cannot look inside the video itself.

At the same time, the individual technologies needed to understand each
of these four channels have matured rapidly. Automatic speech recognition
systems can now transcribe speech with accuracy comparable to a human
transcriber. Optical character recognition can extract printed text from
natural images in real time. Deep-learning object detectors achieve high
precision on standard benchmarks. Large language models can generate
rich textual descriptions of visual scenes. What is missing is a system
that combines all four into a single, end-to-end search pipeline.

This gap motivates the VidSeek project. The goal is to turn any uploaded
video into a fully searchable document by combining all available information
channels and making the result available through a simple natural-language
query interface.

\section{The Essential Question}

The central problem this project addresses is: \textit{how can a user find
a specific moment inside a video using only a natural-language text query,
without watching any part of the video in advance?}

This decomposes into several smaller questions. How do we extract every
piece of searchable information from a video automatically? How do we
represent pieces of information from different modalities - text, audio,
and images - in a common mathematical space so that a single query can
match any of them? How do we store millions of such representations
efficiently so that a query returns results in under a second even on a
large collection? And how do we present the results in a way that is useful
to the user, telling them not just which video contains the answer but
exactly when in the video to look?

VidSeek is designed to answer all of these questions in a single, integrated
system.

\section{Project Objectives}

The overall goal is to build a working multi-modal video search engine that
lets users find specific moments inside a video collection using natural-language
queries. The main objectives are:

\begin{itemize}
  \item \textbf{Multi-modal indexing pipeline}: Build an automated pipeline
    that processes a video through four parallel modules - OCR, ASR, object
    detection, and visual relationship detection - and stores all extracted
    information in a searchable index.

  \item \textbf{Unified embedding space}: Represent text from all four
    modalities as dense vectors in the same embedding space so that a single
    query can retrieve matches from any channel.

  \item \textbf{Efficient vector index}: Design and implement a custom
    two-level Inverted File (IVF) index that can store and search millions
    of embedding vectors with sub-second query latency.

  \item \textbf{Timestamp-level retrieval}: Return search results that
    include precise start and end timestamps, allowing the user to jump
    directly to the relevant moment in the video.

  \item \textbf{REST API and user interface}: Expose the system through a
    clean web API and a browser-based front end so that non-technical users
    can upload videos and run queries without any programming knowledge.
\end{itemize}

\section{Problem Definition and Project Outcomes}

The core problem is one of cross-modal information retrieval. A video is an
unstructured, multi-channel data stream. A user's query is a short piece of
natural-language text. Bridging the gap between these two requires three
things: understanding what is inside the video, encoding that understanding
in a form that can be compared mathematically to the query, and returning
the comparison results in a useful format.

The project is scoped to videos that contain English speech and English
on-screen text, though the architecture does not prevent extension to other
languages. Queries are short natural-language sentences or phrases. The
system is designed for videos of up to two hours in length.

The expected outcomes are:

\begin{itemize}
  \item \textbf{Indexing pipeline}: A software pipeline that takes a video
    file as input and produces a fully indexed, searchable representation.

  \item \textbf{Custom vector store}: A two-level IVF index, implemented
    from scratch in Python, capable of efficient approximate nearest-neighbour
    search over large embedding collections.

  \item \textbf{Relational database schema}: A PostgreSQL schema that stores
    structured metadata - detected objects, visual relationships, OCR regions,
    and transcript segments - alongside the vector index.

  \item \textbf{Search API}: A FastAPI web service that accepts a natural-language
    query and returns a ranked, timestamp-annotated list of matching video
    segments.

  \item \textbf{Web front end}: A browser-based interface for uploading
    videos and querying the system.
\end{itemize}

\section{Document Organization}

The remainder of this report is organised as follows. Chapter~2 presents a
market visibility study, identifying the target users, surveying existing
video search tools, and describing the intended business model.
Chapter~3 reviews the technical background needed to follow the rest of the
report, covering scene segmentation, optical character recognition, object
detection, visual relationship detection, automatic speech recognition, sentence
embeddings, and approximate nearest-neighbour indexing, and then compares a
set of previous works related to video retrieval and multi-modal search.
Chapter~4 describes the overall system design and architecture, including the
input assumptions, the block-level data flow of the indexing and retrieval
pipelines, and a detailed decomposition of each processing module.
Chapter~5 reports how the system was tested and verified, defining the
evaluation metrics and the procedure used to measure retrieval quality.
Chapter~6 concludes the report and outlines directions for future work.
Supporting material - the development platforms and tools, the system use
cases, and a user guide - is collected in the appendices.


% ─────────────────────────────────────────────────────────────────────────────
\chapter{Market Visibility Study}
% ─────────────────────────────────────────────────────────────────────────────

\section{Target Audience}

The fundamental goal of this project is to make video content as searchable
as text documents. The users who would benefit from this are spread across
many fields and can be divided into primary users and key stakeholders.

The primary target audience consists of anyone who needs to find specific
information inside a large collection of videos without watching them
manually. This includes:

\begin{itemize}
  \item \textbf{Students and Researchers}: University students who want to
    find a specific explanation inside a lecture recording, or researchers
    who need to locate a particular finding in a recorded conference talk.
    Manually scrubbing through hours of video is slow and frustrating. A
    text query that returns the exact timestamp saves significant time.

  \item \textbf{Journalists and Editors}: Journalists who manage archives
    of broadcast footage need to locate specific clips quickly. A system
    that can search spoken words, on-screen text, and visible objects in
    archived footage dramatically speeds up their workflow.

  \item \textbf{Legal and Compliance Teams}: Organisations that are required
    by law or company policy to review recorded meetings or security footage
    need a way to find specific events without watching every recording in full.

  \item \textbf{E-learning Platforms}: Online education platforms that host
    hundreds of hours of recorded course content can improve the student
    experience by letting them search within a course for the exact
    explanation of a concept.
\end{itemize}

The secondary audience includes organisations that generate large volumes
of video content and need infrastructure to manage it:

\begin{itemize}
  \item \textbf{Media Companies}: Broadcasters and streaming platforms that
    maintain large archives of recorded programmes.

  \item \textbf{Enterprise Training Departments}: Companies that record
    onboarding sessions, safety briefings, and product demonstrations for
    employees to review later.

  \item \textbf{Healthcare Institutions}: Hospitals and clinics that record
    surgical training sessions or patient consultations for later review
    and education.
\end{itemize}

The scale of the opportunity is significant. According to Cisco, video traffic
accounted for more than 82\% of all internet traffic in 2022, and the volume
continues to grow. The global video surveillance market alone was valued at
over \$45 billion in 2023, a segment where searchable video is extremely valuable.

\section{Market Survey}

The market for video search and retrieval tools exists but is fragmented, and
most existing solutions have important limitations.

\textbf{Platform-level search} (YouTube, Vimeo): The dominant video platforms
offer search over video titles, descriptions, and auto-generated captions.
They do not expose an API for uploading private content and searching it
at the moment level. Caption quality is often poor for technical or domain-specific
vocabulary. There is no support for searching visual content such as detected
objects.

\textbf{Enterprise video platforms} (Kaltura, Panopto, MediaSite): These
products are aimed at university and enterprise customers and include
chapter-level navigation and caption search. They are expensive, require
vendor lock-in, and do not support rich multi-modal queries. Object-level
or relationship-level visual search is not available.

\textbf{AI-powered video analysis services} (Google Video Intelligence API,
AWS Rekognition Video, Azure Video Indexer): These cloud services can
transcribe audio, detect objects, and extract on-screen text from a video.
However, they are pay-per-use, return results in their own proprietary formats,
and do not provide an end-to-end search experience. A developer must build
the indexing and retrieval layer on top of these APIs. Azure Video Indexer
is the closest to a complete solution, but it runs only on Azure infrastructure
and its vector retrieval layer is not customisable.

\textbf{Academic systems}: Several research prototypes support multi-modal
video retrieval, but they are not packaged as deployable products. Most
require domain-specific training data and do not generalise to arbitrary
video content.

The gap that VidSeek addresses is the absence of an open, self-hosted,
end-to-end system that combines all four content modalities (speech, OCR,
objects, relationships) into a single search pipeline, uses a custom vector
index to keep query latency low, and provides a user-friendly web interface.

\section{Business Model}

To ensure sustainability and wide adoption, the proposed business model
combines a free self-hosted option with paid cloud and enterprise tiers.

\textbf{Open-source core}: The core indexing and retrieval engine is released
as open-source software under a permissive licence. This drives adoption,
builds a developer community, and allows academic institutions and individuals
to use the system at no cost.

\textbf{Cloud-hosted SaaS}: A managed cloud version is offered on a
subscription basis. Pricing is tiered by the total number of hours of video
indexed per month. This tier targets small and medium businesses that
want the benefits of VidSeek without managing infrastructure.

\textbf{Enterprise contracts}: Large organisations - broadcasters, universities,
legal firms - require dedicated deployment, custom integrations, and
service-level agreements. These are served through annual enterprise contracts
that include private cloud or on-premise deployment, priority support, and
custom model fine-tuning for domain-specific vocabulary.

\textbf{API access}: A metered API lets developers integrate VidSeek's indexing
and search capabilities into their own applications without deploying the
full stack. Pricing is per video-hour indexed and per query answered.


% ─────────────────────────────────────────────────────────────────────────────
\chapter{Literature Survey}
% ─────────────────────────────────────────────────────────────────────────────

\section{Technical Background Knowledge}

\subsection{Scene Segmentation}

Scene segmentation is the process of dividing a video into a sequence of
non-overlapping clips, where each clip represents a single coherent visual
scene. The boundary between two scenes is marked by an abrupt or gradual
change in visual content, called a shot transition.

The simplest approach compares pairs of consecutive frames using a pixel-level
difference metric. If the difference exceeds a threshold, a cut is detected.
More robust methods compute colour histograms or edge histograms per frame
and compare those instead, making the detector tolerant to small amounts of
motion within a scene.

Deep learning approaches train a binary classifier on pairs of frames, learning
to distinguish a real cut from a fast motion or a flash of light. Systems
such as PySceneDetect and TransNetV2 achieve high precision and recall
on standard benchmarks.

Scene segmentation is the first step in VidSeek because processing every
frame individually would be computationally wasteful. By grouping frames
into scenes, the system can select a representative set of frames from each
scene and process only those.

\subsection{Optical Character Recognition}

Optical Character Recognition (OCR) is the technology that reads text
appearing inside an image. In the context of video, OCR extracts text
from individual frames - subtitles burned into the video, title cards,
lower-third graphics, slide content, or any other text that appears on screen.

A typical OCR pipeline has two stages. The first stage is text detection:
a model draws bounding boxes around all regions of the image that contain
text. The second stage is text recognition: a separate model reads the
characters inside each bounding box and produces a string.

Modern text detection methods use Maximally Stable Extremal Regions (MSER)
to find stable character-candidate blobs, or deep convolutional networks such
as EAST (Efficient and Accurate Scene Text detector) that predict both a
heatmap of text regions and the orientation of the text.

Text recognition is commonly handled by a sequence model. The region inside
the bounding box is fed to a CNN that extracts spatial features, and then
a recurrent decoder (often a CTC-trained LSTM) reads the character sequence
from left to right.

In VidSeek, OCR is applied to a representative frame from each scene. The
extracted text strings are stored in an inverted index and also converted
to embeddings for vector search.

\subsection{Histogram of Oriented Gradients and Support Vector Machines}

Before deep learning became dominant, the most widely used combination for
image classification and object detection was Histogram of Oriented Gradients
(HOG) features fed to a Support Vector Machine (SVM) classifier.

HOG works by dividing an image patch into small cells. For each cell it
computes the magnitude and direction of intensity gradients at every pixel,
and then builds a histogram that counts how many pixels point in each direction.
The histograms of all cells are concatenated into a single feature vector
that describes the local shape of the patch.

An SVM then learns a decision boundary that separates feature vectors that
belong to a given class (for example, a word versus a background region).
SVM with a radial basis function kernel can model non-linear boundaries by
implicitly mapping the features to a higher-dimensional space.

In VidSeek, HOG features and SVM classifiers are used inside the OCR pipeline
to classify candidate text regions and reject non-text noise before passing
regions to the character recognition stage.

\subsection{Object Detection with Faster R-CNN}

Object detection is the task of finding every instance of a set of named
object classes in an image and drawing a bounding box around each one.

Faster R-CNN \cite{ren2015faster} is a two-stage detector that has become
a standard baseline for this task. It consists of two components:

\begin{itemize}
  \item \textbf{Region Proposal Network (RPN)}: A small convolutional network
    that slides over the feature map produced by a backbone CNN and, at each
    position, proposes a fixed set of candidate bounding boxes (called anchor
    boxes) together with an objectness score. The RPN is trained to distinguish
    foreground (any object) from background.

  \item \textbf{RoI Head}: The candidate boxes proposed by the RPN are projected
    back onto the feature map. A fixed-size feature is extracted from each
    region using RoI pooling (or RoI Align). This feature is then passed
    to a fully connected classifier that predicts the object class and refines
    the bounding box coordinates.
\end{itemize}

A Feature Pyramid Network (FPN) is typically added as the backbone to produce
feature maps at multiple scales, allowing the detector to find both small and
large objects in the same image.

In VidSeek, Faster R-CNN is applied to key frames extracted from each scene.
The set of object labels detected in a frame - for example \{``person'', ``car'',
``laptop''\} - is stored in the database and also converted to a text description
that can be embedded and searched.

\subsection{Visual Relationship Detection}

Object detection answers the question ``what objects are in this image?'' but
it does not say anything about how those objects relate to each other.
Visual Relationship Detection (VRD) goes further by identifying triplets of
the form \textit{(subject, predicate, object)}, for example \textit{(person,
holding, phone)} or \textit{(car, parked in front of, building)}.

Early VRD approaches trained a model specifically for this task on datasets
such as Visual Relationship Dataset (VRD) or Visual Genome. The model takes
a pair of detected object bounding boxes as input and classifies the predicate
that connects them.

More recent approaches use large vision-language models or large language
models (LLMs) with image understanding capabilities. Given an image, the
model is prompted to describe the relationships between visible entities in
natural language. This approach generalises better to open-vocabulary predicates
that were not seen during training.

In VidSeek, VRD is performed using the Gemini API. The system sends a key
frame to Gemini with a prompt asking for a structured description of the
relationships between visible objects. The response is parsed and stored as
searchable text.

\subsection{Automatic Speech Recognition}

Automatic Speech Recognition (ASR) is the task of converting a spoken audio
signal into a written transcript.

A classical ASR pipeline consists of three components working together.
An acoustic model converts short windows of audio features into a probability
distribution over phonemes. A pronunciation lexicon maps sequences of phonemes
to words. A language model assigns probabilities to word sequences, favouring
sequences that are grammatically and semantically plausible. A beam-search
decoder combines all three scores to find the most likely word sequence.

Modern end-to-end ASR systems such as Whisper \cite{radford2023whisper}
replace this pipeline with a single large transformer model trained on
hundreds of thousands of hours of transcribed audio. The model takes a
log-mel spectrogram of the audio as input and directly outputs a token sequence.
Whisper is trained on a multilingual, multitask dataset and achieves near-human
word error rates on many languages without any domain-specific fine-tuning.

In VidSeek, audio is extracted from the uploaded video and fed to the
\texttt{openai/whisper-large-v3} model. The output is a time-stamped transcript
where each word has an associated start and end time. The transcript is then
segmented into coherent sentence groups for storage and retrieval.

\subsection{Sentence Embeddings and Transformers}

A sentence embedding maps a piece of text to a fixed-length dense vector
in a high-dimensional space. The key property is that texts with similar
meaning are mapped to vectors that are close together in that space, while
texts with different meanings are mapped to vectors that are far apart.

Early word embedding models such as Word2Vec and GloVe produced per-word
vectors and averaged them to represent a sentence. This was a rough approximation
because it ignored word order and context.

Transformer-based models such as BERT introduced context-sensitive token
representations. A sentence embedding can be derived by taking the output
of the \texttt{[CLS]} token or by averaging all token outputs. Models trained
with a contrastive objective specifically for semantic similarity tasks -
such as Sentence-BERT (SBERT) - produce much better sentence embeddings
than general BERT models.

In VidSeek, a Sentence-BERT style model is used to embed every piece of
text extracted from the video - OCR strings, ASR transcript segments, object
label lists, and VRD descriptions - into the same 768-dimensional vector
space. A user's query is encoded by the same model, and retrieval is performed
by nearest-neighbour search in that space.

\subsection{Approximate Nearest-Neighbour Search with IVF Indexing}

Nearest-neighbour search finds the $k$ vectors in a database that are most
similar to a query vector. The straightforward way to do this is to compute
the distance between the query and every stored vector. This is called
exact nearest-neighbour search and is guaranteed to return the true top-$k$,
but its cost grows linearly with the number of stored vectors, making it
too slow for large collections.

Approximate nearest-neighbour (ANN) search trades a small amount of accuracy
for a large reduction in query time. The most widely adopted family of ANN
algorithms is Inverted File (IVF) indexing, introduced in the context of
image retrieval.

In an IVF index, the vector space is partitioned into $C$ clusters using
k-means. Each cluster is represented by its centroid. At query time, the
system computes the distance from the query to every centroid, selects the
$p$ nearest centroids (where $p \ll C$), and then performs exact search
only within the postings lists of those $p$ clusters. This limits the search
to a small fraction of the total database, reducing query time by a factor
of roughly $C/p$.

VidSeek implements a two-level IVF index. The first level partitions the
global vector space into $C_1$ coarse clusters. Each coarse cluster is
further divided into at most $C_2$ fine sub-clusters. At insertion time,
a vector is first assigned to the nearest level-1 centroid and then to the
nearest level-2 sub-centroid within that group. At query time, $p_1$ level-1
clusters are probed, and within each of them $p_2$ level-2 sub-clusters
are probed, before performing exact search on the small set of candidate
vectors that remain.

This two-level design reduces memory fragmentation and keeps postings lists
small without requiring a full re-clustering of the index when new vectors
are inserted.

\section{Comparative Study of Previous Work}

\subsection{``Video Google: A Text Retrieval Approach to Object Matching in
Videos'' \cite{sivic2003video}}

\subsubsection{Main Problem and Motivation}

This paper introduced the idea of treating video retrieval as an analogue
of text retrieval. The problem it addresses is how to find, in a large
video archive, all shots that contain a particular object without any
annotation.

The motivation was that text retrieval techniques such as TF-IDF and inverted
indexes were already mature and highly efficient. Could the same ideas be
applied to visual features?

\subsubsection{Methodology}

The system extracts local invariant visual features (SIFT descriptors) from
each frame. It then quantises the feature space into a visual vocabulary
(a ``visual dictionary'') using k-means, analogous to the word vocabulary
in text retrieval. Each frame is represented as a bag of visual words, and
an inverted index maps each visual word to the frames that contain it. A
query is specified by selecting an object in a query frame; the system finds
all frames that share many of the same visual words.

\subsubsection{Strengths}

The approach is efficient and scales to large archives. It generalises
naturally from text retrieval, which means mature indexing infrastructure
can be reused. It works without any training on the target object category.

\subsubsection{Weaknesses and Limitations}

The bag-of-words representation discards spatial relationships between features.
It works only for repeated appearance of the exact same object and struggles
with large viewpoint or illumination changes. It does not support natural-language
queries because queries must be specified by example image patches, not text.

\subsection{``Dense Video Captions'' \cite{krishna2017dense}}

\subsubsection{Main Problem and Motivation}

Most video description systems generate a single sentence describing the
whole video. This paper addresses the harder task of generating a separate
natural-language description for every event in a video, together with
temporal localisation of each event.

The motivation is that a video typically contains many distinct events, and
a single-sentence description loses almost all of the content. Fine-grained,
temporally grounded captions enable much more precise retrieval.

\subsubsection{Methodology}

The system uses an activity proposal network to suggest candidate event
intervals in the video. A captioning network then generates a sentence
for each proposed interval, conditioned on both the visual content of
the interval and the captions already generated for other intervals in
the same video (so the descriptions are coherent with each other).
The system is trained end-to-end on the ActivityNet Captions dataset.

\subsubsection{Results}

On the ActivityNet Captions dataset, the system achieved a METEOR score
of 7.29 for joint event detection and captioning, significantly outperforming
earlier single-event captioning baselines.

\subsubsection{Strengths and Limitations}

The dense captioning approach produces rich, searchable text descriptions
of video content. However, training requires a large dataset of manually
annotated event boundaries and captions, which are expensive to collect.
The quality of captions degrades for domain-specific content (medical,
technical) that differs from the training distribution.

\subsection{``CLIP: Learning Transferable Visual Models From Natural Language
Supervision'' \cite{radford2021clip}}

\subsubsection{Main Problem and Motivation}

The problem is how to train an image understanding model that can be applied
to arbitrary new tasks without any task-specific fine-tuning. The standard
approach of training a model on a fixed set of labelled categories is limiting
because the label set is fixed at training time.

\subsubsection{Methodology}

CLIP is trained on 400 million (image, text) pairs collected from the internet.
The training objective is contrastive: an image encoder and a text encoder are
trained so that the embedding of an image and the embedding of its paired
caption are close together in the shared embedding space, while the embeddings
of mismatched pairs are pushed apart.

At inference time, a new category can be introduced simply by writing a natural-language
description of it. The image is embedded by the image encoder, and the
candidate descriptions are embedded by the text encoder. The category whose
text embedding is closest to the image embedding is chosen.

\subsubsection{Strengths}

CLIP generalises to arbitrary visual concepts without retraining. Its shared
embedding space makes it natural to use as a retrieval backbone: queries
written in text can be matched against image or video embeddings directly.

\subsubsection{Limitations}

CLIP does not produce bounding boxes and cannot localise objects spatially.
Its performance on fine-grained or domain-specific visual tasks (for example,
reading text, recognising medical images) lags behind specialised models.
The training data is scraped from the internet and contains biases that
can affect model behaviour.

\subsection{``Whisper: Robust Speech Recognition via Large-Scale Weak
Supervision'' \cite{radford2023whisper}}

\subsubsection{Main Problem and Motivation}

The problem is that ASR models trained on clean speech data degrade sharply
when applied to noisy, accented, or domain-specific audio. Collecting
large amounts of clean, labelled speech data is expensive.

\subsubsection{Methodology}

Whisper uses a standard transformer encoder-decoder architecture. The key
contribution is the training data: 680,000 hours of audio paired with
transcripts, collected from the internet with no manual quality filtering.
The model is trained simultaneously on transcription, translation (from
non-English audio to English text), and voice activity detection.

\subsubsection{Results}

On the LibriSpeech benchmark, Whisper Large V2 achieves 2.7\% word error
rate without any fine-tuning, approaching the performance of models that
were trained specifically on that benchmark.

\subsubsection{Strengths and Limitations}

The zero-shot performance of Whisper across many languages and domains
is a major advance. The model generalises to audio from lectures, interviews,
podcasts, and noisy environments without any task-specific adaptation.
The main limitation is that it is a large model (the large-v3 variant has
1.5 billion parameters) that requires a GPU for real-time inference. It
also sometimes hallucinates text on segments of silence.

\subsection{``Sentence-BERT: Sentence Embeddings using Siamese BERT-Networks''
\cite{reimers2019sentencebert}}

\subsubsection{Main Problem and Motivation}

BERT produces contextualised token embeddings, but deriving a fixed-length
sentence embedding from them is not straightforward. Averaging all token
outputs or using the \texttt{[CLS]} token representation performs poorly
on semantic similarity tasks.

\subsubsection{Methodology}

Sentence-BERT (SBERT) fine-tunes BERT using a Siamese network structure.
Two BERT encoders with shared weights process two input sentences, and the
resulting sentence embeddings are compared using cosine similarity. The
model is trained on the Stanford Natural Language Inference (SNLI) dataset
and the Multi-Genre NLI dataset, which provide pairs of sentences labelled
as entailment, contradiction, or neutral.

\subsubsection{Results}

SBERT achieves a Spearman correlation of 0.869 on the STS-Benchmark
semantic textual similarity task, compared to 0.255 for averaging BERT
token embeddings. Inference time for a sentence pair is 9 ms instead of
65 ms for raw BERT because the sentences are encoded independently.

\subsubsection{Strengths and Limitations}

SBERT embeddings are fast to compute and can be cached in advance, making
them ideal for a retrieval system. The model's limitations are that it is
trained on general-domain text and may not represent highly specialised
vocabulary well without fine-tuning.

\subsection{``Billion-Scale Similarity Search with GPUs'' \cite{johnson2019billion}}

\subsubsection{Main Problem and Motivation}

The problem is how to perform nearest-neighbour search over datasets
containing billions of vectors efficiently on modern GPU hardware.

\subsubsection{Methodology}

The paper describes the FAISS library (Facebook AI Similarity Search).
FAISS implements multiple ANN indexing strategies. The most important for
large-scale use is IVF with product quantisation (IVF-PQ): vectors are
first quantised to reduce memory usage, and then stored in an inverted
file index so that query time is proportional to the fraction of the
database that is probed.

The paper shows that with GPU parallelism, FAISS can query a billion-vector
index with a latency of a few milliseconds.

\subsubsection{Strengths and Limitations}

FAISS is an industrial-strength library used in production at Meta and many
other organisations. It supports both CPU and GPU, and provides exact and
approximate search modes. Its limitation is that it is a C++ library with
a Python wrapper and requires careful memory management. Product quantisation
introduces a recall penalty that can be significant for small databases or
very high precision requirements.


% ─────────────────────────────────────────────────────────────────────────────
\chapter{System Design and Architecture}
% ─────────────────────────────────────────────────────────────────────────────

\section{Overview and Assumptions}

VidSeek is designed as a pipeline system with two distinct phases. The
\textit{indexing phase} processes a video once and stores all extracted
information. The \textit{retrieval phase} accepts a user's query at any
time and searches the stored index.

\subsection{System Input Assumptions}

\begin{itemize}
  \item \textbf{Video format}: The input is a video file in a standard
    container format (MP4, AVI, MOV) that can be decoded by FFmpeg.
  \item \textbf{Language}: Speech and on-screen text are assumed to be
    in English. The ASR module uses Whisper's translate task, which converts
    any language to English, so non-English speech is also supported.
  \item \textbf{Audio quality}: The system assumes that a speech track is
    present. Videos with no audio are indexed only by visual content.
  \item \textbf{Duration}: Videos are assumed to be up to two hours long.
    Longer videos can be processed by splitting them into chunks.
  \item \textbf{Resolution}: The system assumes a minimum resolution of
    480p. Higher resolution improves OCR quality but increases processing time.
  \item \textbf{Query}: The user's query is a natural-language phrase or
    sentence in English.
\end{itemize}

\section{System Architecture}

\subsection{High-Level Block Diagram}

The architecture of VidSeek has two top-level pipelines: the indexing
pipeline and the retrieval pipeline. Figure~\ref{fig:arch_overview}
illustrates the relationship between them.

\begin{figure}[H]
  \centering
  % Replace with actual figure when available
  \fbox{\parbox{0.85\textwidth}{\centering\vspace{3cm}
    [Figure: High-level architecture diagram showing\\
     Indexing Pipeline and Retrieval Pipeline]
  \vspace{3cm}}}
  \caption{High-level architecture of VidSeek.}
  \label{fig:arch_overview}
\end{figure}

\subsection{Indexing Pipeline}

The indexing pipeline is triggered once per video. It runs five processing
modules in a defined order. Figure~\ref{fig:indexing_pipeline} shows the
data flow.

\begin{figure}[H]
  \centering
  \fbox{\parbox{0.85\textwidth}{\centering\vspace{2.5cm}
    [Figure: Indexing pipeline block diagram:\\
     Video $\rightarrow$ Scene Segmenter $\rightarrow$
     \{OCR, Object Detector, VRD\} (parallel) + ASR $\rightarrow$
     Sentence Segmenter $\rightarrow$ Transformer (Embedding) $\rightarrow$
     \{IVF Vector Store, PostgreSQL\}]
  \vspace{2.5cm}}}
  \caption{Indexing pipeline data flow.}
  \label{fig:indexing_pipeline}
\end{figure}

\begin{enumerate}
  \item \textbf{Scene Segmentation}: The video is divided into scenes.
    A representative frame is selected from each scene.
  \item \textbf{Visual Processing} (applied to representative frames):
    \begin{itemize}
      \item OCR extracts on-screen text.
      \item Object Detection detects named objects.
      \item Visual Relationship Detection describes object relationships.
    \end{itemize}
  \item \textbf{Audio Processing}: ASR transcribes the speech in the full
    video with word-level timestamps.
  \item \textbf{Sentence Segmentation}: The word-level transcript is grouped
    into semantically coherent sentence segments.
  \item \textbf{Embedding and Storage}: Every piece of extracted text is
    embedded using a sentence transformer and stored in the IVF vector store.
    Structured metadata is stored in PostgreSQL.
\end{enumerate}

\subsection{Retrieval Pipeline}

The retrieval pipeline is triggered at query time. It is fast because
all heavy processing was done during indexing.

\begin{figure}[H]
  \centering
  \fbox{\parbox{0.85\textwidth}{\centering\vspace{2.5cm}
    [Figure: Retrieval pipeline block diagram:\\
     User Query $\rightarrow$ Sentence Transformer $\rightarrow$
     IVF Search $\rightarrow$ Metadata Lookup (PostgreSQL) $\rightarrow$
     Ranked Results with Timestamps]
  \vspace{2.5cm}}}
  \caption{Retrieval pipeline data flow.}
  \label{fig:retrieval_pipeline}
\end{figure}

\begin{enumerate}
  \item The user's query is encoded by the same sentence transformer used
    during indexing, producing a query vector $\mathbf{q} \in \mathbb{R}^d$.
  \item The IVF index is probed to find the $k$ most similar vectors to
    $\mathbf{q}$.
  \item The metadata associated with each of the $k$ results is fetched
    from PostgreSQL: video path, start timestamp, end timestamp, and the
    original text snippet.
  \item Results are grouped by video and ranked by similarity score.
    The API returns a JSON response with one entry per matching segment,
    including a precise timestamp that can be used to seek to the relevant
    moment.
\end{enumerate}

\section{Time Plan}

Table~\ref{tab:timeplan} shows the project schedule.

\begin{table}[H]
  \centering
  \caption{Project time plan.}
  \label{tab:timeplan}
  \begin{tabular}{llll}
    \toprule
    \textbf{Phase} & \textbf{Task} & \textbf{Start} & \textbf{End} \\
    \midrule
    Phase 1 & Literature review and system design      & Oct 2025 & Nov 2025 \\
    Phase 2 & Scene segmentation and OCR modules       & Nov 2025 & Dec 2025 \\
    Phase 3 & Object detection and VRD modules         & Dec 2025 & Jan 2026 \\
    Phase 4 & ASR and sentence segmentation modules    & Jan 2026 & Feb 2026 \\
    Phase 5 & IVF vector store implementation          & Feb 2026 & Mar 2026 \\
    Phase 6 & Integration, API, and front end          & Mar 2026 & Apr 2026 \\
    Phase 7 & Testing, evaluation, and documentation   & Apr 2026 & Jun 2026 \\
    \bottomrule
  \end{tabular}
\end{table}

\section{Module Descriptions}

\subsection{Scene Segmenter}

\subsubsection{Functional Description}

The Scene Segmenter takes a video file path as input and returns a list
of scenes. Each scene is described by its start frame, end frame, and
the index of a representative frame chosen from the middle of the scene.

\subsubsection{Modular Decomposition}

\paragraph{Frame Reader.}
The Frame Reader decodes the video using FFmpeg and yields frames at a
configurable sampling rate. It attaches a timestamp in seconds to each
frame.

\paragraph{Shot Boundary Detector.}
The Shot Boundary Detector computes a frame-to-frame difference score for
each consecutive pair of frames. When the score crosses a threshold, a
scene boundary is recorded.

\paragraph{Scene Merger.}
Very short scenes (fewer than a configurable minimum number of frames)
are merged with an adjacent scene to avoid fragmenting the video into
segments that are too small to carry meaningful content.

\paragraph{Representative Frame Selector.}
One frame is selected from each scene. This is the frame that is passed
to the OCR, Object Detector, and VRD modules.

\subsubsection{Design Constraints}

The segmenter must process video faster than real time on a CPU to avoid
becoming a bottleneck. The threshold for shot boundary detection is
configurable because different video types (rapid cuts in news footage
versus slow dissolves in lectures) require different settings.

\subsection{OCR Module}

\subsubsection{Functional Description}

The OCR module takes the list of representative frames produced by the
Scene Segmenter and returns a mapping from each scene to the text strings
found in that scene's representative frame.

\subsubsection{Modular Decomposition}

\paragraph{Text Region Detector.}
The Text Region Detector applies MSER to find candidate text blobs in the
frame, then groups nearby blobs into word and line regions. An HOG+SVM
classifier filters out non-text regions.

\paragraph{Character Recogniser.}
Each detected text region is cropped, normalised to a standard height, and
passed to a sequence recognition model. The model outputs the character
sequence as a Unicode string.

\paragraph{Inverted Index Builder.}
The extracted strings are tokenised and used to build an inverted index that
maps each word to the set of scenes in which it appears. This index supports
exact keyword lookup as a complement to semantic vector search.

\subsection{Object Detection Module}

\subsubsection{Functional Description}

The Object Detection module takes the list of representative frames and
returns, for each frame, a list of detected objects. Each detected object
is described by its class label, bounding box coordinates, and confidence
score.

\subsubsection{Modular Decomposition}

\paragraph{Backbone and FPN.}
A ResNet-50 backbone pretrained on ImageNet extracts a rich feature map
from each frame. A Feature Pyramid Network builds feature maps at four
different scales from the backbone outputs, allowing the detector to find
objects of different sizes.

\paragraph{Region Proposal Network.}
The RPN slides a small convolutional network over each scale of the feature
pyramid and proposes candidate bounding boxes together with objectness
scores. Non-maximum suppression removes overlapping proposals.

\paragraph{RoI Classifier.}
For each surviving proposal, a fixed-size feature is extracted from the
FPN feature map using RoI Align. Two fully connected layers then predict
the final class label (from the 91 COCO object categories) and refine
the bounding box.

\paragraph{Result Formatter.}
Detections below a confidence threshold are discarded. The remaining
detections are formatted as a list of (label, box, score) tuples and
saved to an inverted index keyed by object label.

\subsection{Visual Relationship Detection Module}

\subsubsection{Functional Description}

The VRD module takes the list of representative frames and returns, for
each frame, a set of relationship triplets of the form
\textit{(subject, predicate, object)}.

\subsubsection{Approach}

VidSeek uses the Gemini vision API for VRD. Each representative frame is
sent to the API along with a structured prompt that asks the model to list
the visual relationships between objects in the image. The model's response
is parsed into a list of triplets. Triplets are concatenated into a natural-language
sentence of the form ``subject is predicate of object'' for embedding.

This approach avoids training a dedicated VRD model from scratch and
generalises to open-vocabulary predicates.

\subsection{ASR Module}

\subsubsection{Functional Description}

The ASR module takes the video file path as input and returns a list of
word-level transcription entries. Each entry contains the word, its start
time, and its end time, all in seconds from the start of the video.

\subsubsection{Modular Decomposition}

\paragraph{Audio Extractor.}
FFmpeg extracts the audio track from the video and converts it to a 16 kHz
mono WAV file, which is the format expected by Whisper.

\paragraph{Whisper Transcriber.}
The \texttt{whisper-large-v3} model is loaded and applied to the extracted
audio. The model is configured to run in translation mode, converting any
non-English speech to English text. Whisper returns segments with start
and end times and word-level timestamps when the \texttt{word\_timestamps}
option is enabled.

\paragraph{Output Formatter.}
The word-level entries are serialised to a JSON file for use by the
Sentence Segmenter in the next pipeline step.

\subsection{Sentence Segmentation Module}

\subsubsection{Functional Description}

The Sentence Segmenter takes the word-level transcript and groups words
into sentence segments. Each segment is a piece of text with an associated
start and end timestamp.

\subsubsection{Modular Decomposition}

\paragraph{Sliding Window Embedder.}
A sliding window moves over the word-level transcript. At each position it
forms a candidate sentence from the words in the window and embeds it using
the sentence transformer.

\paragraph{Similarity-Based Boundary Detector.}
The cosine similarity between adjacent window embeddings is computed. A
low similarity value indicates a topic change and triggers a segment
boundary. A configurable threshold of 0.75 is used by default.

\paragraph{Segment Aggregator.}
Words between consecutive boundaries are merged into a single segment.
The segment's start time is the start time of its first word, and its end
time is the end time of its last word.

\subsection{Transformer (Embedding) Module}

\subsubsection{Functional Description}

The Transformer module embeds every piece of extracted text into the
shared vector space. It takes the OCR inverted index and the list of
transcript segments as inputs, and produces a list of (vector, metadata)
pairs ready for storage.

\subsubsection{Modular Decomposition}

\paragraph{Text Collector.}
The Text Collector iterates over all OCR entries and transcript segments.
For each entry it attaches a type label (``ocr'' or ``transcript''), the
source video path, and the start and end timestamps.

\paragraph{Batch Encoder.}
Texts are processed in batches using the sentence transformer model
(\texttt{all-MiniLM-L6-v2} by default). Batching is important to use
the GPU or CPU efficiently. The output is a matrix of embeddings of
shape $N \times d$, where $N$ is the number of texts in the batch and
$d$ is the embedding dimension (384 for \texttt{all-MiniLM-L6-v2}).

\paragraph{Storage Interface.}
Each embedding and its associated metadata are passed to the CustomVectorStore
for insertion into the IVF index.

\subsection{Custom Vector Store}

\subsubsection{Functional Description}

The CustomVectorStore is the persistent index that stores all embeddings
and supports approximate nearest-neighbour queries. It wraps the two-level
IVF index implementation.

\subsubsection{Two-Level IVF Index Design}

The index organises vectors in a two-level hierarchy:

\begin{itemize}
  \item \textbf{Level 1} contains up to $C_1 = 500$ coarse cluster centroids.
    Each centroid summarises a broad region of the embedding space.
  \item \textbf{Level 2} contains up to $C_2 = 20$ fine sub-cluster
    centroids per level-1 cluster. Each sub-centroid summarises a narrow
    sub-region within its parent cluster.
\end{itemize}

\paragraph{Insertion.}
When a new vector $\mathbf{v}$ is inserted:
\begin{enumerate}
  \item The $p_1 = 1$ nearest level-1 centroid is found by linear scan.
  \item If $\mathbf{v}$'s distance to the nearest centroid exceeds $\delta_1$
    and the number of level-1 clusters has not reached $C_1$, a new centroid
    is created at $\mathbf{v}$. Otherwise the nearest centroid is used.
  \item The same process is repeated at level 2 within the chosen level-1 cluster.
  \item The vector and its metadata are appended to the postings list of
    the chosen level-2 sub-cluster.
\end{enumerate}

\paragraph{Query.}
When a query vector $\mathbf{q}$ arrives:
\begin{enumerate}
  \item The $p_1 = 12$ nearest level-1 centroids are found.
  \item For each of those $p_1$ clusters, the $p_2 = 6$ nearest level-2
    sub-centroids are found.
  \item The full postings lists of the selected $p_1 \times p_2$ sub-clusters
    are loaded and searched exactly by cosine similarity.
  \item The top-$k$ results across all probed sub-clusters are returned.
\end{enumerate}

\subsubsection{Design Constraints}

The index must support concurrent insertions during indexing and concurrent
reads during retrieval. Thread safety is achieved with per-cluster read-write
locks. The index is stored on disk in binary files for efficiency; the
centroid arrays and metadata are stored in JSON for human readability and
crash recovery.

\subsection{Relational Database Schema}

PostgreSQL is used to store structured metadata. The schema has four main
tables: \texttt{videos}, \texttt{ocr\_entries}, \texttt{objects},
and \texttt{vrd\_entries}. Each table references the parent video by a
foreign key. ASR transcript segments are stored in the vector store's
metadata JSON rather than in a separate table, because they do not require
relational queries.

\subsection{REST API}

The REST API is built with FastAPI and exposes the following endpoints:

\begin{itemize}
  \item \textbf{POST /upload}: Accepts a video file, saves it to disk,
    registers it in the database, and starts the indexing pipeline in a
    background thread. Returns a job ID.
  \item \textbf{GET /status/\{job\_id\}}: Returns the current status of
    an indexing job (pending, running, done, or error) and a progress
    message.
  \item \textbf{POST /search}: Accepts a JSON body with a \texttt{query}
    string and optional \texttt{top\_k} parameter. Returns a ranked list
    of matching segments grouped by video.
  \item \textbf{GET /video/\{video\_name\}}: Streams the video file so
    that the front end can play it and seek to a timestamp.
\end{itemize}

\section{Design Constraints and Trade-offs}

\textbf{Processing time vs. coverage.} Running all four visual modules
on every frame of a long video would take too long. Scene segmentation
solves this by selecting one representative frame per scene, reducing
the number of frames processed by a factor of 10 to 100 depending on
the video.

\textbf{Memory vs. recall.} The two-level IVF index uses a fraction of
the memory that would be required to load all vectors at once for an
exact search, at the cost of a small reduction in recall. The probe
parameters $p_1$ and $p_2$ control this trade-off: increasing them
improves recall but increases query time.

\textbf{Model size vs. speed.} Whisper Large V3 produces better transcripts
than smaller Whisper models, but it requires a GPU and takes longer to run.
For deployments without a GPU, a smaller model (Whisper Base or Small) can
be substituted with a reduction in accuracy.

\textbf{API cost vs. quality for VRD.} Using the Gemini API for VRD avoids
the need to train a dedicated VRD model, but it introduces a per-frame
API cost and a network dependency. An alternative is to use an open-source
vision-language model served locally.


% =============================================================
%  References
% =============================================================
\begin{thebibliography}{99}

\bibitem{sivic2003video}
J.~Sivic and A.~Zisserman,
``Video Google: A text retrieval approach to object matching in videos,''
in \textit{Proceedings of the IEEE International Conference on Computer Vision (ICCV)},
2003, pp.~1470--1477.

\bibitem{krishna2017dense}
R.~Krishna, K.~Hata, F.~Ren, L.~Fei-Fei, and J.~C.~Niebles,
``Dense-captioning events in videos,''
in \textit{Proceedings of the IEEE International Conference on Computer Vision (ICCV)},
2017, pp.~706--715.

\bibitem{radford2021clip}
A.~Radford, J.~W.~Kim, C.~Hallacy, A.~Ramesh, G.~Goh, S.~Agarwal,
G.~Sastry, A.~Askell, P.~Mishkin, J.~Clark, G.~Krueger, and I.~Sutskever,
``Learning transferable visual models from natural language supervision,''
in \textit{Proceedings of the 38th International Conference on Machine Learning (ICML)},
2021, pp.~8748--8763.

\bibitem{radford2023whisper}
A.~Radford, J.~W.~Kim, T.~Xu, G.~Brockman, C.~McLeavey, and I.~Sutskever,
``Robust speech recognition via large-scale weak supervision,''
in \textit{Proceedings of the 40th International Conference on Machine Learning (ICML)},
2023, pp.~28492--28518.

\bibitem{reimers2019sentencebert}
N.~Reimers and I.~Gurevych,
``Sentence-BERT: Sentence embeddings using Siamese BERT-networks,''
in \textit{Proceedings of the 2019 Conference on Empirical Methods in Natural Language Processing (EMNLP)},
2019, pp.~3982--3992.

\bibitem{johnson2019billion}
J.~Johnson, M.~Douze, and H.~Jégou,
``Billion-scale similarity search with GPUs,''
\textit{IEEE Transactions on Big Data},
vol.~7, no.~3, pp.~535--547, 2021.

\bibitem{ren2015faster}
S.~Ren, K.~He, R.~Girshick, and J.~Sun,
``Faster R-CNN: Towards real-time object detection with region proposal networks,''
in \textit{Advances in Neural Information Processing Systems (NeurIPS)},
vol.~28, 2015.

\bibitem{vaswani2017attention}
A.~Vaswani, N.~Shazeer, N.~Parmar, J.~Uszkoreit, L.~Jones, A.~N.~Gomez,
\L.~Kaiser, and I.~Polosukhin,
``Attention is all you need,''
in \textit{Advances in Neural Information Processing Systems (NeurIPS)},
vol.~30, 2017.

\end{thebibliography}

\end{document}
